\scanpage{19}
% archive foliation: 18
$L$ est de torsion dans $\operatorname{Pic}$, $L\otimes\unsure{\underline{\omega}|\omega}^{-1}$ est anti-ample. Donc \uline{tout auto} \deleted{de $P$} \uline{de $P(P^1_{X/\mathbb C}(\unsure{\underline{\omega}|\omega}))$ laisse fixe la section marquée $\sigma$}.

\deleted{OK. Donc sur un $\unsure{\xi|\zeta}$-$P^1(\mathbb C)$ fibré sur $X$, $\exists$ au plus une sect. proj. non stationnaire}
\deleted{\unsure{Application}}
\deleted{\diagram{\begin{tikzcd}
\unsure{M_g^{\xi}(X)|M_{\mathrm{conn}}(X)}\arrow[r,leftrightarrow]&M_g(X)
\end{tikzcd}}}
\deleted{\inserted{$M_{\mathrm{conn}}(X)$}}
\margin{\unsure{Si Si}}

Ceci montre aussi \unsure{en passant} que $E\xrightarrow{\mathrm{épi}}L$ induit par passage au quotient $\mathcal O_X\xrightarrow{\mathrm{épi}}L$, donc $L\simeq\mathcal O_X$, donc on a en fait un autom.\inserted{$u$} de $E$, avec $u(\unsure{\underline{\omega}|\omega})\subset\unsure{\underline{\omega}|\omega}$, et induit \deleted{la}\inserted{l'id} \unsure{sur} le quotient $E/\unsure{\underline{\omega}|\omega}\simeq\mathcal O_X$. Si cet \unsure{autom.} induit $\lambda\in H^0(X,\mathcal O_X^\times)\simeq\mathbb C^\times$ sur $\unsure{\underline{\omega}|\omega}$, \deleted{l'extension} $c\in H^1(X,\unsure{\underline{\omega}|\omega})$ est la classe de l'ex. $E$ de $\mathcal O_X$ par $\unsure{\underline{\omega}|\omega}$, on doit avoir $\lambda c=c$ i.e. $(\lambda-1)c=0$, et comme $c\ne0$, on doit avoir $\lambda=1$. Donc $u$ est un autom. d'extension, donc défini par un hom $\mathcal O_X\longrightarrow\unsure{\underline{\omega}|\omega}$, i.e. par une forme diff $\omega\in H^0(X,\unsure{\underline{\omega}|\omega})$. Donc
\[
\operatorname{Aut}_X(P_X)\inserted{=\operatorname{Aut}_X(P_X,\sigma)}\simeq H^0(X,\unsure{\underline{\omega}|\omega}_X)
\]
\[
u_{\unsure{\underline{\omega}|\omega}}\longleftarrow\!\shortmid\ \unsure{\underline{\omega}|\omega}
\]
(iso. de groupes) La question est alors la suivante : Peut-on avoir $u_{\unsure{\underline{\omega}|\omega}}(c)=\unsure{c'|c}$ si $c$, \deleted{$c'$ \unsure{sont des}}\inserted{est la} conn. sur $P_X$ \unsure{associées} à \deleted{deux}\inserted{une} \uline{conn. projectives} $C$, \deleted{$C'$} ?
